\documentclass[10pt]{extarticle}
\usepackage[a4paper,margin=1.5cm]{geometry}
\usepackage[T1]{fontenc}
\usepackage{amsmath,amssymb}
\usepackage{graphicx}
\usepackage[export]{adjustbox}
\usepackage{multicol}
\usepackage{enumitem}
\usepackage{booktabs,array,tabularx}
\setlength{\columnsep}{0.7cm}
\setlength{\parindent}{0pt}
\setlength{\parskip}{3pt}
\pagestyle{empty}
\makeatletter
\newcommand{\qfigdef}[4]{\@namedef{qf@file@#1}{#2}\@namedef{qf@w@#1}{#3}\@namedef{qf@wide@#1}{#4}}
\newcommand{\qf@check}[1]{\@ifundefined{qf@file@#1}{\PackageError{score2}{Unknown figure id #1}{}}{}}
\newcommand{\qf@img}[4][]{\includegraphics[width=#2\dimexpr\@nameuse{qf@w@#4}pt\relax,max width=#3,max height=0.45\textheight,#1]{\@nameuse{qf@file@#4}}}
\newif\ifqf@full
\newif\ifqf@boxed
\newsavebox{\qf@box}
\newcommand{\qf@open}{\global\qf@fulltrue\par\global\qf@boxedtrue
  \setbox\qf@box\vbox\bgroup\hsize\textwidth\linewidth\textwidth}
\newcommand{\qf@emit}[1]{\par\egroup\global\qf@boxedfalse
  \ifdim\dimexpr\ht\qf@box+\dp\qf@box\relax<#1\textheight\noindent\usebox{\qf@box}\par\else\unvbox\qf@box\fi}
\newcommand{\qfullsplit}{\ifqf@boxed\qf@emit{0.85}\fi}
\newcommand{\qf@close}{\ifqf@boxed\qf@emit{0.45}\fi\global\qf@fullfalse}
\newcommand{\qfullbegin}{\end{multicols}\qf@open}
\newcommand{\qfullend}{\qf@close\begin{multicols}{2}}
\newcommand{\qfullstart}{\qf@open}
\newcommand{\qfulldone}{\qf@close}
\newcommand{\qfig}[2][1]{\qf@check{#2}%
  \ifqf@full\qf@blk{#1}{#2}\else
  \ifnum\@nameuse{qf@wide@#2}=1 \expandafter\@firstoftwo\else\expandafter\@secondoftwo\fi
  {\end{multicols}\begin{center}\qf@img{#1}{\textwidth}{#2}\end{center}\begin{multicols}{2}}%
  {\qf@blk{#1}{#2}}\fi}
\newcommand{\qf@blk}[2]{\par\smallskip{\centering\qf@img{#1}{\linewidth}{#2}\par}\smallskip}
\newcommand{\qfigi}[2][1]{\qf@check{#2}\qf@img[valign=c]{#1}{\linewidth}{#2}}
\makeatother
\newcommand{\papersec}[1]{\par\bigskip{\centering\Large\bfseries PAPER-#1\par}\medskip}
\newcommand{\partsec}[2]{\par\medskip\noindent{\large\bfseries PART-#1: #2}\par\smallskip}
\newcommand{\sect}[1]{\par\smallskip\noindent\textbf{#1}\par}
\newcommand{\qn}[1]{\par\smallskip\noindent\textbf{#1.}~}
\newcommand{\paraq}[2]{\par\smallskip\noindent\textbf{Paragraph for Questions #1 and #2}\par}
\newenvironment{qoptions}{\begin{itemize}[leftmargin=1.8em,itemsep=1pt,topsep=2pt,label={}]}{\end{itemize}}
\newcommand{\opt}[2]{\item[(#1)] #2}
% figure registry (generated by score2_tools.py finalize)
\qfigdef{p1_phy_s1_q01_1}{JEE Advanced (2021 Series) - Test 07 (25 Apr 2021)-img/p1_phy_s1_q01_1.png}{139.0}{0}
\qfigdef{p1_phy_s1_q02_1}{JEE Advanced (2021 Series) - Test 07 (25 Apr 2021)-img/p1_phy_s1_q02_1.png}{89.9}{0}
\qfigdef{p1_phy_s1_q03_1}{JEE Advanced (2021 Series) - Test 07 (25 Apr 2021)-img/p1_phy_s1_q03_1.png}{81.4}{0}
\qfigdef{p1_phy_s1_q05_1}{JEE Advanced (2021 Series) - Test 07 (25 Apr 2021)-img/p1_phy_s1_q05_1.png}{203.5}{0}
\qfigdef{p1_phy_s1_q05_A}{JEE Advanced (2021 Series) - Test 07 (25 Apr 2021)-img/p1_phy_s1_q05_A.png}{170.0}{0}
\qfigdef{p1_phy_s1_q05_B}{JEE Advanced (2021 Series) - Test 07 (25 Apr 2021)-img/p1_phy_s1_q05_B.png}{177.8}{0}
\qfigdef{p1_phy_s1_q05_C}{JEE Advanced (2021 Series) - Test 07 (25 Apr 2021)-img/p1_phy_s1_q05_C.png}{184.2}{0}
\qfigdef{p1_phy_s1_q05_D}{JEE Advanced (2021 Series) - Test 07 (25 Apr 2021)-img/p1_phy_s1_q05_D.png}{191.5}{0}
\qfigdef{p1_phy_s2_q07_1}{JEE Advanced (2021 Series) - Test 07 (25 Apr 2021)-img/p1_phy_s2_q07_1.png}{102.5}{0}
\qfigdef{p1_phy_s2_q08_1}{JEE Advanced (2021 Series) - Test 07 (25 Apr 2021)-img/p1_phy_s2_q08_1.png}{146.8}{0}
\qfigdef{p1_phy_s2_q09_1}{JEE Advanced (2021 Series) - Test 07 (25 Apr 2021)-img/p1_phy_s2_q09_1.png}{140.5}{0}
\qfigdef{p1_phy_s2_q10_1}{JEE Advanced (2021 Series) - Test 07 (25 Apr 2021)-img/p1_phy_s2_q10_1.png}{174.7}{0}
\qfigdef{p1_phy_s2_q11_2}{JEE Advanced (2021 Series) - Test 07 (25 Apr 2021)-img/p1_phy_s2_q11_2.png}{137.3}{0}
\qfigdef{p1_phy_s3_q02_2}{JEE Advanced (2021 Series) - Test 07 (25 Apr 2021)-img/p1_phy_s3_q02_2.png}{88.2}{0}
\qfigdef{p1_phy_s3_q03_1}{JEE Advanced (2021 Series) - Test 07 (25 Apr 2021)-img/p1_phy_s3_q03_1.png}{110.4}{0}
\qfigdef{p1_phy_s3_q05_1}{JEE Advanced (2021 Series) - Test 07 (25 Apr 2021)-img/p1_phy_s3_q05_1.png}{167.8}{0}
\qfigdef{p1_phy_s3_q06_1}{JEE Advanced (2021 Series) - Test 07 (25 Apr 2021)-img/p1_phy_s3_q06_1.png}{245.2}{0}
\qfigdef{p1_chem_s1_q01_1}{JEE Advanced (2021 Series) - Test 07 (25 Apr 2021)-img/p1_chem_s1_q01_1.png}{362.2}{1}
\qfigdef{p1_chem_s1_q02_1}{JEE Advanced (2021 Series) - Test 07 (25 Apr 2021)-img/p1_chem_s1_q02_1.png}{371.1}{1}
\qfigdef{p1_chem_s1_q05_A}{JEE Advanced (2021 Series) - Test 07 (25 Apr 2021)-img/p1_chem_s1_q05_A.png}{87.7}{0}
\qfigdef{p1_chem_s1_q05_B}{JEE Advanced (2021 Series) - Test 07 (25 Apr 2021)-img/p1_chem_s1_q05_B.png}{87.7}{0}
\qfigdef{p1_chem_s1_q05_C}{JEE Advanced (2021 Series) - Test 07 (25 Apr 2021)-img/p1_chem_s1_q05_C.png}{87.7}{0}
\qfigdef{p1_chem_s1_q05_D}{JEE Advanced (2021 Series) - Test 07 (25 Apr 2021)-img/p1_chem_s1_q05_D.png}{87.7}{0}
\qfigdef{p1_chem_s1_q06_A}{JEE Advanced (2021 Series) - Test 07 (25 Apr 2021)-img/p1_chem_s1_q06_A.png}{118.9}{0}
\qfigdef{p1_chem_s1_q06_B}{JEE Advanced (2021 Series) - Test 07 (25 Apr 2021)-img/p1_chem_s1_q06_B.png}{114.7}{0}
\qfigdef{p1_chem_s1_q06_C}{JEE Advanced (2021 Series) - Test 07 (25 Apr 2021)-img/p1_chem_s1_q06_C.png}{104.9}{0}
\qfigdef{p1_chem_s1_q06_D}{JEE Advanced (2021 Series) - Test 07 (25 Apr 2021)-img/p1_chem_s1_q06_D.png}{110.0}{0}
\qfigdef{p1_chem_s2_q09_A}{JEE Advanced (2021 Series) - Test 07 (25 Apr 2021)-img/p1_chem_s2_q09_A.png}{67.5}{0}
\qfigdef{p1_chem_s2_q09_B}{JEE Advanced (2021 Series) - Test 07 (25 Apr 2021)-img/p1_chem_s2_q09_B.png}{66.7}{0}
\qfigdef{p1_chem_s2_q09_C}{JEE Advanced (2021 Series) - Test 07 (25 Apr 2021)-img/p1_chem_s2_q09_C.png}{66.5}{0}
\qfigdef{p1_chem_s2_q09_D}{JEE Advanced (2021 Series) - Test 07 (25 Apr 2021)-img/p1_chem_s2_q09_D.png}{66.5}{0}
\qfigdef{p1_chem_s2_q10_2}{JEE Advanced (2021 Series) - Test 07 (25 Apr 2021)-img/p1_chem_s2_q10_2.png}{74.8}{0}
\qfigdef{p1_chem_s2_q10_3}{JEE Advanced (2021 Series) - Test 07 (25 Apr 2021)-img/p1_chem_s2_q10_3.png}{75.1}{0}
\qfigdef{p1_chem_s2_q10_1}{JEE Advanced (2021 Series) - Test 07 (25 Apr 2021)-img/p1_chem_s2_q10_1.png}{46.9}{0}
\qfigdef{p1_chem_s3_q05_1}{JEE Advanced (2021 Series) - Test 07 (25 Apr 2021)-img/p1_chem_s3_q05_1.png}{154.9}{0}
\qfigdef{p1_chem_s3_q06_1}{JEE Advanced (2021 Series) - Test 07 (25 Apr 2021)-img/p1_chem_s3_q06_1.png}{398.0}{1}
\qfigdef{p2_phy_s1_q01_1}{JEE Advanced (2021 Series) - Test 07 (25 Apr 2021)-img/p2_phy_s1_q01_1.png}{115.8}{0}
\qfigdef{p2_phy_s1_q02_1}{JEE Advanced (2021 Series) - Test 07 (25 Apr 2021)-img/p2_phy_s1_q02_1.png}{110.4}{0}
\qfigdef{p2_phy_s1_q03_1}{JEE Advanced (2021 Series) - Test 07 (25 Apr 2021)-img/p2_phy_s1_q03_1.png}{162.1}{0}
\qfigdef{p2_phy_s1_q04_1}{JEE Advanced (2021 Series) - Test 07 (25 Apr 2021)-img/p2_phy_s1_q04_1.png}{206.8}{0}
\qfigdef{p2_phy_s1_q05_1}{JEE Advanced (2021 Series) - Test 07 (25 Apr 2021)-img/p2_phy_s1_q05_1.png}{171.9}{0}
\qfigdef{p2_phy_s1_q06_1}{JEE Advanced (2021 Series) - Test 07 (25 Apr 2021)-img/p2_phy_s1_q06_1.png}{203.8}{0}
\qfigdef{p2_phy_s1_q08_1}{JEE Advanced (2021 Series) - Test 07 (25 Apr 2021)-img/p2_phy_s1_q08_1.png}{212.9}{0}
\qfigdef{p2_phy_s2_q11_1}{JEE Advanced (2021 Series) - Test 07 (25 Apr 2021)-img/p2_phy_s2_q11_1.png}{85.5}{0}
\qfigdef{p2_phy_s2_q11_3}{JEE Advanced (2021 Series) - Test 07 (25 Apr 2021)-img/p2_phy_s2_q11_3.png}{85.5}{0}
\qfigdef{p2_phy_s2_q11_5}{JEE Advanced (2021 Series) - Test 07 (25 Apr 2021)-img/p2_phy_s2_q11_5.png}{85.5}{0}
\qfigdef{p2_phy_s2_q11_7}{JEE Advanced (2021 Series) - Test 07 (25 Apr 2021)-img/p2_phy_s2_q11_7.png}{85.5}{0}
\qfigdef{p2_phy_s2_q11_2}{JEE Advanced (2021 Series) - Test 07 (25 Apr 2021)-img/p2_phy_s2_q11_2.png}{74.6}{0}
\qfigdef{p2_phy_s2_q11_4}{JEE Advanced (2021 Series) - Test 07 (25 Apr 2021)-img/p2_phy_s2_q11_4.png}{72.6}{0}
\qfigdef{p2_phy_s2_q11_6}{JEE Advanced (2021 Series) - Test 07 (25 Apr 2021)-img/p2_phy_s2_q11_6.png}{73.1}{0}
\qfigdef{p2_phy_s2_q11_8}{JEE Advanced (2021 Series) - Test 07 (25 Apr 2021)-img/p2_phy_s2_q11_8.png}{72.0}{0}
\qfigdef{p2_phy_s2_q12_1}{JEE Advanced (2021 Series) - Test 07 (25 Apr 2021)-img/p2_phy_s2_q12_1.png}{75.3}{0}
\qfigdef{p2_phy_s3_q01_1}{JEE Advanced (2021 Series) - Test 07 (25 Apr 2021)-img/p2_phy_s3_q01_1.png}{139.8}{0}
\qfigdef{p2_phy_s3_q02_1}{JEE Advanced (2021 Series) - Test 07 (25 Apr 2021)-img/p2_phy_s3_q02_1.png}{140.9}{0}
\qfigdef{p2_phy_s3_q06_1}{JEE Advanced (2021 Series) - Test 07 (25 Apr 2021)-img/p2_phy_s3_q06_1.png}{128.1}{0}
\qfigdef{p2_phy_s3_q08_1}{JEE Advanced (2021 Series) - Test 07 (25 Apr 2021)-img/p2_phy_s3_q08_1.png}{143.1}{0}
\qfigdef{p2_chem_s1_q03_1}{JEE Advanced (2021 Series) - Test 07 (25 Apr 2021)-img/p2_chem_s1_q03_1.png}{84.0}{0}
\qfigdef{p2_chem_s1_q03_2}{JEE Advanced (2021 Series) - Test 07 (25 Apr 2021)-img/p2_chem_s1_q03_2.png}{70.9}{0}
\qfigdef{p2_chem_s1_q03_3}{JEE Advanced (2021 Series) - Test 07 (25 Apr 2021)-img/p2_chem_s1_q03_3.png}{70.9}{0}
\qfigdef{p2_chem_s1_q03_4}{JEE Advanced (2021 Series) - Test 07 (25 Apr 2021)-img/p2_chem_s1_q03_4.png}{194.2}{0}
\qfigdef{p2_chem_s1_q08_1}{JEE Advanced (2021 Series) - Test 07 (25 Apr 2021)-img/p2_chem_s1_q08_1.png}{171.5}{0}
\qfigdef{p2_chem_s1_q08_2}{JEE Advanced (2021 Series) - Test 07 (25 Apr 2021)-img/p2_chem_s1_q08_2.png}{117.4}{0}
\qfigdef{p2_chem_s1_q08_3}{JEE Advanced (2021 Series) - Test 07 (25 Apr 2021)-img/p2_chem_s1_q08_3.png}{117.4}{0}
\qfigdef{p2_chem_s2_q11_1}{JEE Advanced (2021 Series) - Test 07 (25 Apr 2021)-img/p2_chem_s2_q11_1.png}{131.1}{0}
\qfigdef{p2_chem_s3_q06_1}{JEE Advanced (2021 Series) - Test 07 (25 Apr 2021)-img/p2_chem_s3_q06_1.png}{177.6}{0}
\qfigdef{p2_math_s2_q12_1}{JEE Advanced (2021 Series) - Test 07 (25 Apr 2021)-img/p2_math_s2_q12_1.png}{108.5}{0}

\begin{document}
\begin{center}{\Large\bfseries JEE Advanced (2021 Series) --- Test 07 (25 Apr 2021)}\end{center}
\begin{multicols}{2}
\papersec{1}
\partsec{1}{PHYSICS}
\sect{Section-I(i)}
\qn{1} Crank $OA$ rotating around point $O$, leads to movement of wheel 1 of radius $R=20$ cm rolling inside circle surface 2. Wheel 1, in contact with wheel 3, makes it spin around points $O$ (Fig.). Find the ratio $\omega_3/\omega$ of angular speed of wheel 3 and that of crank $OA$ if the radius of wheel 3 is $r=10$ cm.
\qfig{p1_phy_s1_q01_1}
\begin{qoptions}
\opt{A}{$1$}
\opt{B}{$2$}
\opt{C}{$3$}
\opt{D}{$6$}
\end{qoptions}
\qn{2} Two small bodies of masses $1$ kg and $4$ kg are at rest on a smooth horizontal surface. When they are at a separation of $10$ m, they are released to move under mutual attractive interaction. Then, which of the following statements is \textbf{CORRECT}?
\qfig{p1_phy_s1_q02_1}
\begin{qoptions}
\opt{A}{They will never meet}
\opt{B}{They will meet at a distance of $6$ m from $A$.}
\opt{C}{To an observer at the centre of mass, the particles have equal and opposite momenta.}
\opt{D}{They will meet at a distance of $8$ m from $B$.}
\end{qoptions}
\qn{3} A hollow thin hemisphere of mass $M$ and radius $R$ lies on a flat horizontal rubber surface. In the upper hemisphere there is a small hole in which a long thin tube is inserted (Fig.). Hemisphere tube is filled with water. To what maximum height can be poured the water into the tube?
\qfig{p1_phy_s1_q03_1}
\begin{qoptions}
\opt{A}{$\dfrac{M}{\pi R^2\rho}+\dfrac23R$}
\opt{B}{$\dfrac{M}{\pi R^2\rho}+\dfrac13R$}
\opt{C}{$\dfrac{M}{2\pi R^2\rho}+\dfrac23R$}
\opt{D}{$\dfrac{M}{\pi R^2\rho}+\dfrac43R$}
\end{qoptions}
\qn{4} A particle of mass $m$ is moving in potential energy $U=\dfrac12kx^2+\dfrac{\alpha}{2x^2}$ where $k$ and $\alpha$ are positive constants. The angular frequency of small oscillations for the simple harmonic motion of the particle about a stable equilibrium position is
\begin{qoptions}
\opt{A}{$\sqrt{\dfrac km}$}
\opt{B}{$\sqrt{\dfrac{2k}{m}}$}
\opt{C}{$2\sqrt{\dfrac km}$}
\opt{D}{$4\sqrt{\dfrac km}$}
\end{qoptions}
\qn{5} A uniform rigid meter-scale is held horizontally with one of its end at the edge of a table and the other supported by hand. Some coins of negligible mass are kept on the meter scale as shown in the figure.
\qfig{p1_phy_s1_q05_1}
As the hand supporting the scale is removed, the scale starts rotating about its edge on the table and the coins start moving. If a photograph of the rotating scale is taken soon after, it will look closest to:
\begin{qoptions}
\opt{A}{\qfigi{p1_phy_s1_q05_A}}
\opt{B}{\qfigi{p1_phy_s1_q05_B}}
\opt{C}{\qfigi{p1_phy_s1_q05_C}}
\opt{D}{\qfigi{p1_phy_s1_q05_D}}
\end{qoptions}
\qn{6} There are three planets in circular orbits around a fixed star at distances $a$, $4a$ and $9a$, respectively in same sense (say anticlockise). At time $t=0$, the star and the three planets are in a straight line. The period of revolution of the closest planet is $T$. What minimum time after $t=0$ will they again be in the same straight line?
\begin{qoptions}
\opt{A}{$8\,T$}
\opt{B}{$27\,T$}
\opt{C}{$108\,T$}
\opt{D}{$512\,T$}
\end{qoptions}
\sect{Section-I(ii)}
\qn{7} In an arrangement shown in the figure, $d\ll D$, where $d$ is the distance of separation of the slits $S_1$ and $S_2$ and $D$ is the distance between the slits and the screen. Source $S$ is emitting monochromatic light of wavelength $\lambda$. Then
\qfig{p1_phy_s2_q07_1}
\begin{qoptions}
\opt{A}{minimum value of $d$ for the dark fringe at $O$ is $\sqrt{\dfrac\lambda2}$.}
\opt{B}{minimum value of $d$ for the dark fringe at $O$ is $\sqrt{\dfrac\lambda3}$.}
\opt{C}{for this value of $d$, the fringe width is $\sqrt{2\lambda}$.}
\opt{D}{for this value of $d$, the fringe width is $\sqrt{\dfrac\lambda2}$.}
\end{qoptions}
\qn{8} A thermally insulated chamber of volume $2V_0$ is divided by a frictionless insulating piston of area $S$ into two equal parts $A$ and $B$. Part $A$ has an ideal gas at pressure $P_0$ and temperature $T_0$ and in part $B$ is vacuum. A massless spring of force constant $K$ is connected with the piston and the wall of the container as shown. Initially the spring is unstretched. The ideal gas in chamber $A$ is allowed to expand. Let in equilibrium the spring be compressed by a length $x_0$, then
\qfig{p1_phy_s2_q08_1}
\begin{qoptions}
\opt{A}{Final pressure of the gas is $Kx_0/S$}
\opt{B}{Work done by the gas is $\tfrac12 kx_0^2$}
\opt{C}{Magnitude of change in internal energy of the gas is $\dfrac12kx_0^2$}
\opt{D}{Temperature of the gas in equilibrium is less than initial temperature}
\end{qoptions}
\qn{9} Two capacitors of capacitance $C_1$ and $C_2$ are charged to a potential difference of $V_1$ and $V_2$ respectively and are connected to an inductor of inductance $L$ as shown in the figure. Initially key $k$ is open. Now key $k$ is closed and current in the circuit starts increasing. When current in the circuit is maximum (Given that $C_1V_1$ is not equals $C_2V_2$)
\qfig{p1_phy_s2_q09_1}
\begin{qoptions}
\opt{A}{charge on both the capacitors will be different.}
\opt{B}{induced emf in the inductor is zero}
\opt{C}{magnitude of potential difference across both the capacitors is same}
\opt{D}{energy stored in the inductor is $\dfrac12\dfrac{C_1C_2}{C_1+C_2}(V_1+V_2)^2$}
\end{qoptions}
\qn{10} Further support for the photon model of electromagnetic waves comes from \textit{Compton scattering,} in which x-rays scatter from electrons, changing direction and frequency in the process. Classical electromagnetic wave theory cannot explain the change in frequency of the x rays on scattering, but the photon model can. Suppose an x-ray photon is moving to the right. It has a collision with a slow-moving electron, as in figure. The photon transfers energy and momentum to the electron, which recoils at a high speed. The x-ray photon loses energy, and the photon energy formula $E=hf$ tells us that its frequency must decrease. The collision looks very much like the collision between two particles.
\qfig{p1_phy_s2_q10_1}
\begin{qoptions}
\opt{A}{When the x-ray photon scatters from the electron, it's speed stays the same.}
\opt{B}{When the x-ray photon scatters from the electron, it's speed decreases.}
\opt{C}{When the electron is struck by the x-ray photon, it's de Broglie wavelength decreases.}
\opt{D}{When the electron is struck by the x-ray photon, it's de Broglie wavelength stays the same.}
\end{qoptions}
\qn{11} An electron at point $A$ in the Figure has a velocity $v_0=10^7\ \mathrm{ms^{-1}}$. Then (if $e=1.6\times10^{-19}$ C, $AB=10$ cm and $m=9.1\times10^{-31}$ kg)
\qfig{p1_phy_s2_q11_2}
\begin{qoptions}
\opt{A}{magnitude of magnetic induction that will cause the electron to follow the semicircular path from $A$ to $B$ is $1.125$ mT}
\opt{B}{time taken to go from $A$ to $B$ is $1.57\times10^{-7}$ s}
\opt{C}{angular momentum of the electron about $A$ when $t=0.785\times10^{-7}$ s is $4.5\times10^{-24}\ \mathrm{kgm^2s^{-1}}$.}
\opt{D}{angle made by the velocity vector with the $x$-axis at $t=\dfrac38T$ is $-\dfrac\pi4$ rad.}
\end{qoptions}
\qn{12} Two identical, stretched wires are vibrated together. They produce few beats per second. $T_1$ and $T_2$ denote the tensions in the two wires at any instant and $T_1>T_2$ (initially). How can the change in $T_1$ or $T_2$ may be performed so that beat frequency increases.
\begin{qoptions}
\opt{A}{$T_1$ decreased}
\opt{B}{$T_1$ increased}
\opt{C}{$T_2$ decreased}
\opt{D}{$T_2$ increased}
\end{qoptions}
\sect{Section-II}
\qn{1} A smooth insulating fixed guide ring of radius $R$ is in horizontal plane. There is a small bead of mass $m$ and carrying charge $q$ and constraint to move along guide. The guide is placed in a coaxially cylindrically symmetric uniform magnetic field of strength $B$ which is perpendicular to the guide plane. Starting from $t=0$, the magnetic field is changed to $B=B_0+\alpha t$. If the force on the bead by the guide at time $t$ is $\dfrac{\alpha q^2Rt}{Xm}(YB_0+\alpha t)$, find $\dfrac{XY}{X+Y}$. Assume gravity to be absent.
\qn{2} The shaded area in the figure is the side view of a disk shaped magnetic field region of radius $R$ and thickness $d$. A parallel electron beam being accelerated by a voltage $V$ is normally incident onto the region. If the spatial distribution of the magnetic field in the region such that the beam is focused to a point along the disk central axis at a distance $L\ (\gg R)$ from the disk is $B=\dfrac{y^\eta}{d^\alpha L^\beta}\sqrt{\dfrac{\lambda mV}{e}}$, find the value of $\dfrac{5\alpha}{\lambda+\beta+\eta}$. Given $m$ and $e$ are the electron's mass and charge and $y$ is the distance of incident electron from axis of disc.
\qfig{p1_phy_s3_q02_2}
\qn{3} Figure shows a uniformly charged spherical shell $B$ of charge ($q_B=q$) kept inside uniformly charged spherical shell $A$ of charge ($q_A=q$). Let electric field due to $A$ and $B$ at any point are $\vec E_A$ and $\vec E_B$ respectively. If $dV$ represents elementary volume, then the value of $\dfrac{\int_{\text{entire space}}\vec E_B\cdot\vec E_B\,dV}{\int_{\text{entire space}}\vec E_A\cdot\vec E_B\,dV}$ will be
\qfig{p1_phy_s3_q03_1}
\qn{4} The zero of main scale matches with zero of vernier scale. $9^{\text{th}}$ division of main scale matches with $10^{\text{th}}$ division of vernier. No other match is present. After sliding, zero of vernier is found to be in between $3$ \& $4^{\text{th}}$ divisions of main scale \& $5^{\text{th}}$ division of vernier matches with $x^{\text{th}}$ division of main scale. The value of $x/3$ is
\qn{5} A wheel with six spokes is positioned perpendicular to a uniform magnetic field $B$ of magnitude $0.5$ T. The field is directed into the plane of the paper and is present over the entire region of the wheel as shown in the figure. When the switch $S$ is closed, there is an initial current of $6$ A between the axle and the rim; and the wheel begins to rotate. The resistance of the spokes and of the rim is negligible. The radius of the wheel is $0.2$ m. Calculate the initial torque (in Nm) on the wheel.
\qfig{p1_phy_s3_q05_1}
\qn{6} The displacement vs time graph for two sound waves $A$ and $B$ which travel along the same medium are shown in the figure. Their loudness are $\beta_A$ and $\beta_B$. If $X=(\beta_A-\beta_B)$, find the value of $X$.
\qfig{p1_phy_s3_q06_1}
\qn{7} An ammeter and a voltmeter are connected in series to a battery with an emf $E=6.0$ V. When a certain resistance is connected in parallel with the voltmeter, the reading of the voltmeter decreases two times, whereas the reading of the ammeter increases the same number of times. Find the voltmeter reading (in volt) after the connection of the resistance.
\qn{8} The image of an object formed by a convex spherical mirror is magnified one third in comparison with the object. If the mirror is displaced by $40$ cm along its principal optical axis in the direction of light, the image is again found one-third in comparison with the object. Find the radius of curvature of the mirror (in cm).
\qfullbegin
\partsec{2}{CHEMISTRY}
\sect{Section-I(i)}
\qn{1} Consider the following compounds :-
\qfig{p1_chem_s1_q01_1}
\qfullsplit
The correct order of reactivity with $\mathrm{Br_2+CH_3COOH}$ is
\begin{qoptions}
\opt{A}{I $<$ II $<$ III $<$ IV}
\opt{B}{I $>$ II $>$ III $>$ IV}
\opt{C}{II $>$ I $>$ III $>$ IV}
\opt{D}{II $>$ I $>$ IV $>$ III}
\end{qoptions}
\qn{2} The compounds those would not respond to tests of both nitrogen and sulfur with sodium fusion extracts are?
\qfig{p1_chem_s1_q02_1}
\qfullsplit
\begin{qoptions}
\opt{A}{I and III}
\opt{B}{III and IV}
\opt{C}{I and IV}
\opt{D}{II and IV}
\end{qoptions}
\qfullend
\qn{3} Which of the following does liberate hydrogen gas when treated with water under suitable condition ?
\par\noindent (I) $\mathrm{B_2H_6}$\quad (II) $\mathrm{B_3N_3H_6}$\quad (III) $\mathrm{SiH_4}$\quad (IV) $\mathrm{XeF_4}$
\begin{qoptions}
\opt{A}{I, III, IV only}
\opt{B}{II, III only}
\opt{C}{I, II, IV only}
\opt{D}{I, II, III only}
\end{qoptions}
\qn{4} Select the correct statement when potassium ferrocyanide is treated with concentrated sulphuric acid ?
\begin{qoptions}
\opt{A}{$\mathrm{Fe^{2+}}$ of $\mathrm{\left[Fe(CN)_6\right]^{4-}}$ is oxidised to $\mathrm{FeO_2^-}$}
\opt{B}{CO gas is evolved}
\opt{C}{HCOOH is one of final product}
\opt{D}{HCN gas is one of the final products}
\end{qoptions}
\qn{5} The graph that represents the temperature (T) -- entropy (S) variation of a carnot cycle is
\begin{qoptions}
\opt{A}{\qfigi{p1_chem_s1_q05_A}}
\opt{B}{\qfigi{p1_chem_s1_q05_B}}
\opt{C}{\qfigi{p1_chem_s1_q05_C}}
\opt{D}{\qfigi{p1_chem_s1_q05_D}}
\end{qoptions}
\qn{6} The Boyle temperature ($T_B$) is defined as the temperature at which the properties of a real gas coincide with those of an ideal gas in the low pressure limit. The graph that shows the pressure dependance of the compressibility factor ($Z$) for a real gas at $T_B$ is
\begin{qoptions}
\opt{A}{\qfigi{p1_chem_s1_q06_A}}
\opt{B}{\qfigi{p1_chem_s1_q06_B}}
\opt{C}{\qfigi{p1_chem_s1_q06_C}}
\opt{D}{\qfigi{p1_chem_s1_q06_D}}
\end{qoptions}
\sect{Section-I(ii)}
\qn{7} Which of the following compounds exist as polymeric solid at room temperature ?
\begin{qoptions}
\opt{A}{$\mathrm{NaPO_3}$}
\opt{B}{$\mathrm{H_3BO_3}$}
\opt{C}{$\mathrm{BeH_2}$}
\opt{D}{$\mathrm{BeCl_2}$}
\end{qoptions}
\qn{8} The reactions which produce $\mathrm{Na_2S_2O_3}$ are
\begin{qoptions}
\opt{A}{$\mathrm{Na_2S+I_2+Na_2SO_3\longrightarrow\,?}$}
\opt{B}{$\mathrm{Na_2S+Na_2CO_3+SO_2\longrightarrow\,?}$}
\opt{C}{$\mathrm{S+NaOH\longrightarrow\,?}$}
\opt{D}{$\mathrm{Na_2SO_3+S\xrightarrow{\ \Delta\ }\,?}$}
\end{qoptions}
\qn{9} Consider the following reaction :
\[\text{(D)}-\text{glucose}\xrightarrow[\text{cat.\,H}^+]{\mathrm{Ph{-}NH{-}NH_2}\ \text{(3 equiv)}}\mathrm X\]
Among the following, the compound(s) whose osazone derivative(s) will have the same melting point as that of X is/are
\begin{qoptions}
\opt{A}{\qfigi{p1_chem_s2_q09_A}}
\opt{B}{\qfigi{p1_chem_s2_q09_B}}
\opt{C}{\qfigi{p1_chem_s2_q09_C}}
\opt{D}{\qfigi{p1_chem_s2_q09_D}}
\end{qoptions}
\qn{10} Consider the following compounds
\par\noindent\qfigi{p1_chem_s2_q10_2}\ I\quad\qfigi{p1_chem_s2_q10_3}\ II\quad\qfigi{p1_chem_s2_q10_1}\ III\par
Select the correct statements
\begin{qoptions}
\opt{A}{(I) undergoes decarboxylation faster than II or III on heating}
\opt{B}{(I) when heated in the presence of acetophenone produces a compound having tertiary alcoholic group}
\opt{C}{Compound (I) acts as bidentate ligand}
\opt{D}{The compound (III) has the highest melting point.}
\end{qoptions}
\qn{11} A complex is repesented as $\mathrm{CoCl_3\cdot xNH_3}$. Its $0.1$ molal solution in water shows $\Delta T_f=0.558^\circ$. If the cryscopic and ebullioscopic constants for $\mathrm{H_2O}$ are $1.86$ K-kg mol$^{-1}$ and $0.52$ K-kg mol$^{-1}$ respectively, identify the correct versions assuming $100\%$ ionisation of the complex.
\begin{qoptions}
\opt{A}{Hexaaminecobalt(III) chloride is the IUPAC name for the complex}
\opt{B}{The above solution will boil between $100^\circ$C to $101^\circ$C}
\opt{C}{The central metal ion is $\mathrm{d^2sp^3}$ hybridised}
\opt{D}{The coordination sphere is incapable of exhibiting geometrical isomerism}
\end{qoptions}
\qn{12} You have a mixture of three amino acids : $A(\mathrm{pI}=3)$, $B(\mathrm{pI}=5)$ and $C(\mathrm{pI}=10)$, where pI represents their isoelectric points. Under electrophoresis at pH $=8$, the movement of each component of the mixture is considered. Which of the following options are incorrect regarding movement acids under electrophoresis :
\begin{qoptions}
\opt{A}{A and B to anode; C to cathode}
\opt{B}{A to anode; B stationary; C to cathode}
\opt{C}{A to cathode ; B stationary; C to anode}
\opt{D}{A to anode; B and C to cathode}
\end{qoptions}
\sect{Section-II}
\qn{1} For $3s$ orbital of hydrogen atom, the normalised wave function is given by
\[\psi_{3s}=\dfrac1{81\sqrt{3\pi}}\left(\dfrac1{a_0}\right)^{3/2}\left[27-\dfrac{18r}{a_0}+\dfrac{2r^2}{a_0^2}\right]e^{-\frac{r}{3a_0}}\]
The above mentioned orbital ($3s$) has two nodes at $2a_0$ and $xa_0$. Find the value of $x/4$.\\
{[}Use : $3\sqrt3=5${]}
\qn{2} The emf of the cell,
\[\mathrm{M(s)\,|\,M_nX_n\,(0.01M)\,||\,M_nX_n\,(0.1M)\,/\,M(s)}\]
is found to be $0.0295$ V at $25^\circ$C. If $\mathrm{M_nX_n}$ is a soluble salt, find the value of $2n\,/\,3$.
\qn{3} $2.616$ gm of an element `X' is heated with NaOH and $\mathrm{NaNO_3}$ to produce $\mathrm{Na_2XO_2}$ and $\mathrm{NH_3}$. Ammonia produced is absorbed in $100$ml of $0.11$ M $\mathrm{H_2SO_4}$ and excess of acid is back titrated with $48$ ml of $0.25$ M NaOH. If $N$ moles of hydrogen gas would be liberated when $2$ moles of `X' is treated with excess of NaOH, find the value of $\dfrac1{2N}$.\\
{[}Atomic mass : Cr $=52$, Zn $=65.4$, Fe $=56$, Cu $=63.5$, Ni $=28${]}
\qn{4} When aqueous solution of sodium phathalate is electrolysed the product (hyrdocarbon) at anode has $x$ degree of unsaturation. Final the value of $\dfrac x2$.
\qn{5} Consider the following reaction given below
\qfig{p1_chem_s3_q05_1}
Find the value of $\dfrac n4$ where $n$ is the total number of methyl groups present in one molecule of (X).
\qfullbegin
\qn{6} Specific rotations of freshly prepared aqueous solutions of \textbf{I} and \textbf{II} are $+112$ and $+18$, respectively. On standing the optical rotation of aqueous solution of \textbf{I} slowly decreases give a final value of $+52.7$ due to equilibration with \textbf{II}. Under this state of equilibrium, what the ratio \textbf{II} : \textbf{I}?
\qfig{p1_chem_s3_q06_1}
\qfullend
\qn{7} Find the value of $\dfrac n5$ where $n$ is the total number of compounds among the following which produce $\mathrm{H_2S}$ gas on treatement with water.
\par\noindent $\mathrm{P_4S_{10}}$, $\mathrm{Al_2S_3}$, MgS, $\mathrm{Na_2S}$, BaS, $\mathrm{Cr_2S_3}$
\qn{8} Consider the following compounds / ions $\mathrm{SnCl_2}$, $\mathrm{HgCl_2}$, $\mathrm{H_2O}$, NOCl, $\mathrm{H_2S}$, $\mathrm{I_3^+}$, $\mathrm{OF_2}$, $\mathrm{NH_2^-}$, $\mathrm{C_2H_2}$. Find the value of $\dfrac x2$ where $x$ is the total number of species having bent shape around central atom.
\partsec{3}{MATHEMATICS}
\sect{Section-I(i)}
\qn{1} A variable plane cuts off intercepts from the co-ordinate axes which are equal to the roots of the equation $x^3+5x=p-qx^2$ ($p,q$ are real numbers). The locus of the foot of the perpendicular from the origin to the plane is
\begin{qoptions}
\opt{A}{$(x^2+y^2+z^2)^2(xy+yz+zx)=5$}
\opt{B}{$(x^2+y^2+z^2)^4\left(\dfrac1{xy}+\dfrac1{yz}+\dfrac1{zx}\right)=5$}
\opt{C}{$(x^2+y^2+z^2)^2\left(\dfrac1{xy}+\dfrac1{yz}+\dfrac1{zx}\right)=5$}
\opt{D}{$(x^2+y^2+z^2)(xy+yz+zx)=5$}
\end{qoptions}
\qn{2} Consider $f(x)=x^4-3x^3+2$ and let $h(x)=\dfrac{f(x)}{g(x)}$ be a rational function such that\\
(i) Domain of $h(x)$ is $x\in\mathbb R-\{-1,1\}$.\\
(ii) $g(x)$ does not have any repeated linear factors.\\
(iii) As $x$ approaches infinity, $h(x)$ also approaches $-\infty$.\\
(iv) $\displaystyle\lim_{x\to1}h(x)=\dfrac12$ then $\displaystyle\lim_{x\to0}\dfrac{g(\cos x)}{5\ell n(\sec^2x)}$ is equal to -
\begin{qoptions}
\opt{A}{$1$}
\opt{B}{$2$}
\opt{C}{$3$}
\opt{D}{$5$}
\end{qoptions}
\qn{3} If $A$ is a square matrix of order $5$ and $9A^{-1}=4A^T$ then number of digit before decimal in $|\operatorname{adj}(\operatorname{adj}(\operatorname{adj}A))|$ (where $A^{-1}$, $A^T$ and $\operatorname{adj}(A)$ denotes the inverse, transpose and adjoint of matrix $A$ respectively) is : ($\log3=0.477$, $\log2=0.303$)
\begin{qoptions}
\opt{A}{$56$ digits}
\opt{B}{$60$ digits}
\opt{C}{$58$ digits}
\opt{D}{$53$ digits}
\end{qoptions}
\qn{4} Let `$d$' be defined as $d=\sqrt{a^2+5a+8}-\sqrt{a^2-3a+4}$, $a\in R$. The number of integral values of $d$ is
\begin{qoptions}
\opt{A}{$7$}
\opt{B}{$9$}
\opt{C}{$13$}
\opt{D}{$15$}
\end{qoptions}
\qn{5} In $\triangle ABC$ with usual notation, if $a,c,A$ are given and $b_1=2b_2$ where $b_1,b_2$ are two values of the third side then the value of $\dfrac{c\sqrt{1+8\sin^2A}}{a}$ is less than
\begin{qoptions}
\opt{A}{$1$}
\opt{B}{$2$}
\opt{C}{$3$}
\opt{D}{$4$}
\end{qoptions}
\qn{6} Sum of the solutions of the equation $\sin^2\theta\tan\theta+\cos^2\theta\cot\theta-\sin2\theta=1+\tan\theta+\cot\theta$, where $\theta\in(0,\pi)-\left\{\dfrac\pi2\right\}$, is -
\begin{qoptions}
\opt{A}{$3\pi$}
\opt{B}{$\dfrac{3\pi}2$}
\opt{C}{$\dfrac{5\pi}2$}
\opt{D}{$7\pi$}
\end{qoptions}
\sect{Section-I(ii)}
\qn{7} If $z$ is a unimodular complex number such that
\[\operatorname{Re}(z-1)+\operatorname{Re}(z^2)=\int_0^{\pi/2}\sin x\,\ln|\sin x-\cos x|\,dx,\]
then $z$ can satisfies
\begin{qoptions}
\opt{A}{$z+\bar z=-2$}
\opt{B}{$z+\bar z=1$}
\opt{C}{$\arg(z)=\dfrac\pi3$}
\opt{D}{$\arg z=\pi$}
\end{qoptions}
\qn{8} Let $f:[0,1]\to\mathbb R$ be a continuous function. If $f(x)$ is differentiable on $(0,1)$ \& $\displaystyle\int_0^1f(x)\,dx=0$, then which option(s) is/are correct :-
\begin{qoptions}
\opt{A}{$\displaystyle\int_0^cf(x)\,dx=f(c)$ for at least one $c$ in $(0,1)$}
\opt{B}{$\displaystyle\int_0^cf(x)\,dx=f(c)$ for all $c\in(0,1)$}
\opt{C}{$\displaystyle\int_0^cf(x)\,dx=-\left(\dfrac{1-c}{c}\right)f(c)$ for at least one $c$ in $(0,1)$}
\opt{D}{$\displaystyle\int_0^cf(x)\,dx=-\left(\dfrac{1-c}{c}\right)f(c)$ for all $c\in(0,1)$}
\end{qoptions}
\qn{9} Consider a curve $C:y=\sqrt{\dfrac{3x^2-21}{7}}$ and a line $L$ of slope $1$, which is tangent to $C$ at point $A$, then
\begin{qoptions}
\opt{A}{Equation of $L$ is $y=x+2$}
\opt{B}{Equation of $L$ is $y=x-2$}
\opt{C}{Equation of circle passing through origin and touches the curve $C$ at $A$ is $x^2+y^2+x-41y=0$}
\opt{D}{Equation of circle passing through origin and touches the curve $C$ at $A$ is $4x^2+4y^2+x-41y=0$}
\end{qoptions}
\qn{10} Let $|a|,|b|<1$, where $a$ and $b$ are complex number, then
\begin{qoptions}
\opt{A}{$|a-b|<|1-\bar ab|$}
\opt{B}{$|a-b|>|1-\bar ab|$}
\opt{C}{$\left|\dfrac{a-b}{1-a\bar b}\right|<1$}
\opt{D}{$\left|\dfrac{a-b}{1-a\bar b}\right|>1$}
\end{qoptions}
\qn{11} Let $f(x)=\dfrac{x^3-4}{(x-1)^3}$ then
\begin{qoptions}
\opt{A}{$\displaystyle\lim_{x\to-\infty}\sqrt{f(x)}=1$}
\opt{B}{$x=-2$ is the point of local minimum.}
\opt{C}{the equation $f(x)=c$ will have three real and distinct roots for $c\in\left(\dfrac49,1\right)\cup(1,4)$.}
\opt{D}{the line $y=1$ intersects the graph of $f(x)$ at exactly one point.}
\end{qoptions}
\qn{12} Given a differentiable function $f(x)$ satisfying the integral equation $f(x)=\displaystyle\int_0^xe^{x-y}f'(y)\,dy-(x^2-x+1)e^x\ \forall x\in\mathbb R$, then-
\begin{qoptions}
\opt{A}{$f(x)$ is local minimum at $x=-\dfrac12$}
\opt{B}{$f(x)$ has exactly one asymptote}
\opt{C}{If $f(x)$ has an extremum point at $x=x_1$ and point of inflection at $x=x_2$, then $|x_1-x_2|=1$}
\opt{D}{$|f'(x)-f(x)|=|f(x)-f''(x)|$}
\end{qoptions}
\sect{Section-II}
\qn{1} Two circles $C_1$ and $C_2$ both passes through the point $A(1,2)$ and $E(2,1)$ and touch the line $4x-2y=9$ at $B$ and $D$ respectively. The possible coordinates of a point $C$ such that the quadrilateral $ABCD$ is a parallelogram is $(a,b)$ then the value of $|ab|$ is :
\qn{2} Given $A(0,0,1)$, $B(2,1,2)$ be two points in 3-D space also $P(\alpha,\beta)$ ($\alpha>0$, $\beta>0$) be a point on the $xy$-plane such that $\angle APB$ is maximum, then $\left[\dfrac{\alpha+3\beta}{5}\right]+\left(\dfrac{\alpha+2}{\beta+1}\right)$ is equal to\\
(where $[\,\cdot\,]$ greatest integer function and $A,P,B$ are non collinear points)
\qn{3} A pack of playing card is consist of $51$ cards. Cards are drawn one by one from the pack without replacement. If first $4$ cards drawn consist of only one honour card, then the probability that the missing card being honour card is $\dfrac ab$ (where $a$ and $b$ are co-prime and all Jacks, Queens, Kings and Aces are considered as honour card) then $\dfrac ab$ is \underline{\hspace{1.5cm}}.
\qn{4} A certain kind of bacteria either die, split into two or split into three bacteria. All splits are exact copies. The chances of dying is $\dfrac14$, the chances of splitting into two is $\dfrac12$ and splitting into three is $\dfrac14$. If the probability that it survives for infinite length of time is $\dfrac{m-\sqrt{13}}{n}$ ($m,n\in\mathbb N$), then the value of $(m/n)$ is
\qn{5} Let $f(x,y)$ is defined for $x,y\in\mathbb N$ and satisfy $f(x,x)=x$
\begin{gather*}
f(x,x+y)=\left(1+\dfrac xy\right)\cdot f(x,y)\\
f(x,y)=f(y,x).
\end{gather*}
Then $f(52,14)$ is
\qn{6} There are five different boxes and seven different balls. The number of ways in which these balls can be distributed so that box $2$ and box $4$ contains only $1$ ball each and at least $1$ box is empty is $N$. (Order of putting the balls in the boxes is NOT considered). Then the value of $N/100$ is \underline{\hspace{1.5cm}}.
\qn{7} Let $f(x)$ and $g(x)$ are differentiable function of $x$ such that
\begin{gather*}x^2\,\bigl(g(f(x))\bigr)\cdot\bigl(f'(g(x))\bigr)\cdot\bigl(g'(x)\bigr)\\ =(1-2x^2)\bigl(g'(f(x))\bigr)\cdot\bigl(f'(x)\bigr)\cdot f(g(x))\ \ \forall,\ x\in R.\end{gather*}
Moreover $g(x)$ is positive and $f(x)$ is non negative and $\displaystyle\int_0^kf(g(x))\,dx=1-e^{-k^2}$, then the value of $\ln\left(\dfrac{g(f(3))}{g(f(2))}\right)$ is \underline{\hspace{1.5cm}}
\qn{8} If the product of all possible values of $c\in\mathbb N$ for which $x^2-20x+c=0$ has always a real root, is $n!$ then sum of digits of $n$, is
\papersec{2}
\partsec{1}{PHYSICS}
\sect{Section-I(i)}
\qn{1} Consider a uniformly charged sheet $ABCD$, which is a part of an equilateral triangular sheet of side $a$ as shown in the figure. Choose the \textbf{CORRECT} options regarding the electric field $E$ at point $O$ due to this sheet.
\qfig{p2_phy_s1_q01_1}
\begin{qoptions}
\opt{A}{Magnitude of $\vec E$, increases with the increase in $a$ (Keeping charge density same)}
\opt{B}{Magnitude of $E$, decreases with increase in $a$ (keeping total charge same)}
\opt{C}{If charge density is $\sigma$ and $a=1$ m, magnitude of $E$ is equal to $\dfrac{7\sigma}{44\varepsilon_0}\ln\sqrt2$.}
\opt{D}{If charge density is $\sigma$ and $a=2$ m, magnitude of $E$ is equal to $\dfrac{7\sigma}{22\varepsilon_0}\ln\sqrt2$.}
\end{qoptions}
\qn{2} A light stick of length $\ell$ rests with its one end against the smooth wall and other end against the smooth horizontal floor as shown in the figure. The bug starts at rest from point $B$ and moves such that the stick always remains at rest, $a_p$ is the magnitude of acceleration of bug of mass $m$, which depends upon its distance of $x$ from the top end of the stick. Choose the \textbf{CORRECT} option(s) :
\qfig{p2_phy_s1_q02_1}
\begin{qoptions}
\opt{A}{$a_P=\dfrac{g}{\sin\theta}\left(1-\dfrac x\ell\right)$}
\opt{B}{$a_P=\dfrac{g}{\cos\theta}\left(1-\dfrac x\ell\right)$}
\opt{C}{The time taken by the bug to reach the bottom of the stick having started at the top end from rest is $\dfrac\pi2\sqrt{\dfrac{\ell\sin\theta}{3g}}$.}
\opt{D}{The time taken by the bug to reach the bottom of the stick having started at the top end from rest is $\dfrac\pi2\sqrt{\dfrac{\ell\sin\theta}{g}}$.}
\end{qoptions}
\qn{3} A very large slab of insulating material having face area $A$ and dielectric constant $K=2$ is placed between plates of parallel plate capacitor and some uniform charges on plates are given. The system is placed in uniform electric field (fig). The field is perpendicular to the surface of the material.
\qfig{p2_phy_s1_q03_1}
\begin{qoptions}
\opt{A}{The total charge on right surface is $Q$}
\opt{B}{The free charge on left surface is $3Q$}
\opt{C}{The bound charge on right surface is $2Q$}
\opt{D}{The magnitude of the applied electric field is $2Q/A\varepsilon_0$.}
\end{qoptions}
\qn{4} Two strings made of same material, same cross-section but different lengths are fixed at both ends and excited in their respective fundamental modes such that total oscillation energy is the same in both cases. Let $f_1,f_2$ be the oscillation frequencies and $A_1,A_2$ be the amplitudes of oscillation of the antinodes in the two cases. Assuming both strings to be under same tension, pick the \textbf{CORRECT} option(s).
\qfig{p2_phy_s1_q04_1}
\begin{qoptions}
\opt{A}{$\dfrac{f_1}{f_2}=\dfrac14$}
\opt{B}{$\dfrac{A_1}{A_2}=\dfrac12$}
\opt{C}{$\dfrac{f_1}{f_2}=\dfrac41$}
\opt{D}{$\dfrac{A_1}{A_2}=\dfrac21$}
\end{qoptions}
\qn{5} A uniform rod of mass $m$ and length $\ell$ lies on a horizontal frictionless table. A particle of same mass moving with a speed $v_0$ sticks to one end of the rod at $t=0$. At $t=t$
\qfig{p2_phy_s1_q05_1}
\begin{qoptions}
\opt{A}{force exerted by particle on rod will be $\dfrac9{25}\dfrac{mv_0^2}{\ell}$}
\opt{B}{just after collision acceleration of point $B$ of rod in $y$-direction is $\dfrac{25}{27}\dfrac{v_0^2}{\ell}$}
\opt{C}{In the frame of centre of mass of system, path of centre of the rod is circle.}
\opt{D}{Path of centre of rod is cycloid.}
\end{qoptions}
\qn{6} A bullet of mass $m$ travelling horizontally with speed $u$ strikes a wooden block of mass `$M$' placed on a smooth horizontal plane. The penetration is assumed uniform and the bullet comes to rest with respect to block after penetrating a distance `$d$' into the block. Then choose the \textbf{CORRECT} statement(s):
\qfig{p2_phy_s1_q06_1}
\begin{qoptions}
\opt{A}{Ultimately the total loss of kinetic energy is $\dfrac{mMu^2}{2(M+m)}$}
\opt{B}{The value of the resistance force (assumed constant) offered by the wood is $F=\dfrac{mMu^2}{2d(M+M)}$}
\opt{C}{The distance covered by the bullet w.r.t. the ground before it comes to rest w.r.t. block is not equal to $\dfrac{dm}{M+m}$}
\opt{D}{The block moves less distance than the bullet wrt ground.}
\end{qoptions}
\qn{7} Hydrogen gas absorbs radiations of wavelength $\lambda_0$ and consequently emit radiations of $6$ different wavelengths of which three wavelengths are shorter than $\lambda_0$. Choose the \textbf{CORRECT} statement(s).
\begin{qoptions}
\opt{A}{The final excited state of the atom is $n=4$}
\opt{B}{The initial state of the atom may be $n=2$}
\opt{C}{There are two transitions in balmer series}
\opt{D}{There are three transitions belonging to Lyman series.}
\end{qoptions}
\qn{8} In the figure shown, resistance of each resistor are same. If current in $AG$ is $5$ A then
\qfig{p2_phy_s1_q08_1}
\begin{qoptions}
\opt{A}{current in $FL$ is $5$ A}
\opt{B}{current in $BA$ is $15$ A}
\opt{C}{Current in $AF$ from top resistor is $10$ A}
\opt{D}{current in $BH$ is $20$ A}
\end{qoptions}
\sect{Section-I(ii)}
\qn{9} A particle is taken to a distance $r\ (>R)$ from centre of earth. $R$ is the radius of earth. It is given velocity $v$ which is perpendicular to $\vec r$ (measured from centre of earth). Match the following assuming particle do not collide with earth. Mass of particle ($m$) is negligible in comparision to mass of earth ($M$) and assume there is no air resistance :
\par\noindent\textbf{List--I}
\begin{qoptions}
\opt{P}{$v=\sqrt{\dfrac{GM}{r}}$}
\opt{Q}{$v=\sqrt{\dfrac{2GM}{r}}$}
\opt{R}{$\sqrt{\dfrac{2GM}{r}}>v>\sqrt{\dfrac{GM}{r}}$}
\opt{S}{$\sqrt{\dfrac{GM}{r}}>v>\sqrt{\dfrac{GM}{r}\left(\dfrac{2R}{r+R}\right)}$}
\end{qoptions}
\par\noindent\textbf{List--II}
\begin{qoptions}
\opt{1}{Distance of point mass from centre of earth first increases then decreases}
\opt{2}{Distance of point mass from centre of earth first decreases then increases}
\opt{3}{Distance of point mass from centre of earth monotonically increases}
\opt{4}{Distance of point mass from centre of earth remain constant}
\end{qoptions}
\begin{qoptions}
\opt{A}{P$\to$4; Q$\to$3; R$\to$1; S$\to$2}
\opt{B}{P$\to$3; Q$\to$4; R$\to$2; S$\to$1}
\opt{C}{P$\to$1; Q$\to$2; R$\to$3; S$\to$4}
\opt{D}{P$\to$4; Q$\to$3; R$\to$2; S$\to$1}
\end{qoptions}
\qn{10} Three ants are sitting on the vertices of an equilateral triangle of side of $a$. At $t=0$, ant 1 starts approaching ant 2 with a speed $v$, ant 2 starts approaching ant 3 with a speed $v$ and ant 3 start approaching ant 1 with a speed $v$. If $a=1$m, $v=1$m/s match the following
\par\noindent\textbf{List--I}
\begin{qoptions}
\opt{P}{The three ant meet at the centroid of the triangle at time (in s)}
\opt{Q}{The angular velocity (in rad/s) of ant 2 with respect to ant 1 at $t=0$ is}
\opt{R}{The angular acceleration (in rad/s$^2$) of ant 2 with respect to ant 1 at $t=0$ is}
\opt{S}{The area of the equilateral triangle (in m$^2$) formed by the three ants at any time $t=a/4v$ will be}
\end{qoptions}
\par\noindent\textbf{List--II}
\begin{qoptions}
\opt{1}{$\dfrac{\sqrt3}2$}
\opt{2}{$\dfrac{25\sqrt3}{256}$}
\opt{3}{$\dfrac23$}
\opt{4}{$\dfrac{3\sqrt3}4$}
\end{qoptions}
\begin{qoptions}
\opt{A}{P$\to$1; Q$\to$4; R$\to$2; S$\to$3}
\opt{B}{P$\to$3; Q$\to$1; R$\to$4; S$\to$2}
\opt{C}{P$\to$4; Q$\to$3; R$\to$1; S$\to$2}
\opt{D}{P$\to$3; Q$\to$1; R$\to$1; S$\to$4}
\end{qoptions}
\qn{11} In a series AC circuit if we change either resistance or reactance not only current amplitude changes, but phase angle between current and voltage also changes. If we take voltage as reference and draw current phasor then head of this phasor moves on a particular trajectory. In List-I, variation of resistance or reactance is specified and in List-II is given path followed by head of current phasor.
\par\noindent\textbf{List--I}
\begin{qoptions}
\opt{P}{\qfigi{p2_phy_s2_q11_1}}
\opt{Q}{\qfigi{p2_phy_s2_q11_3}}
\opt{R}{\qfigi{p2_phy_s2_q11_5}}
\opt{S}{\qfigi{p2_phy_s2_q11_7}}
\end{qoptions}
\par\noindent\textbf{List--II}
\begin{qoptions}
\opt{1}{\qfigi{p2_phy_s2_q11_2}}
\opt{2}{\qfigi{p2_phy_s2_q11_4}}
\opt{3}{\qfigi{p2_phy_s2_q11_6}}
\opt{4}{\qfigi{p2_phy_s2_q11_8}}
\end{qoptions}
\begin{qoptions}
\opt{A}{P$\to$1; Q$\to$3; R$\to$2; S$\to$4}
\opt{B}{P$\to$4; Q$\to$1; R$\to$3; S$\to$2}
\opt{C}{P$\to$2; Q$\to$3; R$\to$1; S$\to$4}
\opt{D}{P$\to$3; Q$\to$2; R$\to$1; S$\to$4}
\end{qoptions}
\qn{12} Figure shows a box $B$ of mass $m$ held fixed at large height $H$ (which is very large in comparison to dimension of the box, but very small in comparison to radius of earth). A particle $P$ whose mass is equal to that of box is suspended at $O$ (which is centre of mass of the box) with a massless string. Particle can swing in vertical plane without colliding box. Particle is pulled aside to horizontal position and released. The box is held fixed till the particle reaches lowest position and then box is released at $t=0$. $t$ is a general time before box collides ground.
\qfig{p2_phy_s2_q12_1}
\par\noindent\textbf{List-I}
\begin{qoptions}
\opt{P}{Tension in the thread at the $t=t$ in the units of $mg$ will be}
\opt{Q}{Speed of box at time $t=\pi\sqrt{\dfrac{\ell}{2g}}$ in the units of $\sqrt{\dfrac{g\ell}{8}}\sqrt{\pi^2+4}$ will be}
\opt{R}{Acceleration of the box at $t=\pi\sqrt{\dfrac{\ell}{2g}}$ in the units of $g$ will be}
\opt{S}{Initial acceleration ($t=0^+$) of particle in the units of $g$ will be}
\end{qoptions}
\par\noindent\textbf{List-II}
\begin{qoptions}
\opt{1}{$1$}
\opt{2}{$2$}
\opt{3}{$0$}
\opt{4}{none}
\end{qoptions}
\begin{qoptions}
\opt{A}{P$\to$3; Q$\to$2; R$\to$1; S$\to$4}
\opt{B}{P$\to$1; Q$\to$1; R$\to$3; S$\to$4}
\opt{C}{P$\to$4; Q$\to$1; R$\to$3; S$\to$2}
\opt{D}{P$\to$1; Q$\to$2; R$\to$3; S$\to$3}
\end{qoptions}
\sect{Section-II}
\qn{1} As shown, a piston chamber of cross section area $A$ is filled with ideal gas. A sealed piston of mass $m$ is right at the middle height of the cylinder at equilibrium. The friction force between the chamber wall and the piston can be ignored. The mass of the rest of the chamber is $M$. The atmosphere pressure is $P_0$. Now slowly pull the piston upwards isothermally, find the maximum value of $\dfrac Mm$ such that the chamber can be lifted off the ground. Given that $P_0=\dfrac{3mg}{2A}$
\qfig{p2_phy_s3_q01_1}
\qn{2} Initially a disc of radius $30$ cm was in rest \& we apply force $F$ as shown in figure. If friction acting on disc is zero then find the value of $x$ (in cm).
\qfig{p2_phy_s3_q02_1}
\qn{3} A glass bulb with volumetric expansion coefficient $\alpha=1.50\times10^{-5}\ {}^0\mathrm C^{-1}$ is weighed in liquid at temperatures $t=10^0$C and $t_1=20^0$C. The weights of the displaced liquids are $p=100.01$ g and $p_1=100$ g respectively. Find the volumetric expansion coefficient $\alpha_1$ (in the units of $10^{-5}\ {}^0\mathrm C^{-1}$) of the liquid in the temperature interval from $t$ to $t_1$.
\qn{4} A mercury drop of radius $1$ cm is sprayed into $10^6$ drops of equal size. Calculate the energy (in mJ) spent for the process (Surface Tension of Mercury $=35\times10^{-3}\ \mathrm{Nm^{-1}}$).
\qn{5} One end of a string of length $L=15$m is tied to the ceiling of an elevator accelerating upwards with an acceleration $20$ m/s$^2$. The other end of the string is free. The linear mass density of the string varies linearly from $0$ to $\lambda_0$ from bottom to top. Find the time taken by a pulse (in seconds) to reach from bottom to the top. ($g=10\ \mathrm{ms^{-2}}$)
\qn{6} A square frame of wire connected to a capacitor is kept in a magnetic field which varies with position as well as time and is given as $\vec B=\alpha xt(-\hat k)$ as shown in figure, where $\alpha$ is a constant. The charge on the capacitor at any time is $\dfrac n2C\alpha\ell^3$, Find $n$.
\qfig{p2_phy_s3_q06_1}
\qn{7} A radioactive sample has decay constant $\lambda$. The rate of production of nuclei in the given sample is $\dfrac{9\lambda N_0^2}{N}$, where $N_0$ is the number of radioactive nuclei in the sample at $t=0$ and $N$ is the number of radioactive nuclei in the sample at time $t=t$ sec. If the number of nuclei present in the radioactive sample at $t\to\infty$ is $m\times10^6$ nuclei then find the value of $m$. (Given $N_0=10^6$ nuclei)
\qn{8} A tube of uniform cross-section is used to siphon water from a vessel, as shown in the Figure. Assuming that the cross-section of the vessel is very much greater than that of the tube, find the maximum value of height $h_1$ (in m), for which the siphon will work. [Take $h_2=3$ m, atmospheric pressure $P_{atm}=10^5\ \mathrm{Nm^{-2}}$, density of water $\rho=1000\ \mathrm{kg\text{-}m^{-3}}$ and acceleration due to gravity $g=10\,\mathrm{ms^{-2}}$]
\qfig{p2_phy_s3_q08_1}
\partsec{2}{CHEMISTRY}
\sect{Section-I(i)}
\qn{1} Which of the following produce(s) gaseous product on treatment with dilute HCl ?
\begin{qoptions}
\opt{A}{$\mathrm{Na_2S_2O_3}$}
\opt{B}{$\mathrm{NaHCO_3}$}
\opt{C}{$\mathrm{Na_2SO_3}$}
\opt{D}{$\mathrm{Na_2S}$}
\end{qoptions}
\qn{2} The complex(s) which can be oxidised on exposure into air or on boiling is / are
\begin{qoptions}
\opt{A}{$\mathrm{\left[Fe(H_2O)_6\right]Cl_2}$}
\opt{B}{$\mathrm{K_4\left[Co(CN)_6\right]}$}
\opt{C}{$\mathrm{Na_4\left[Co(NO_2)_6\right]}$}
\opt{D}{$\mathrm{\left[Co(NH_3)_6\right]Cl_3}$}
\end{qoptions}
\qn{3} Consider the reaction given below :
\par\noindent\qfigi{p2_chem_s1_q03_1}$\ \xrightarrow{\mathrm{PCl_5}}\ \underset{\text{major}}{(\mathrm x)}\ \xrightarrow{\mathrm{NH_3}}\ \underset{\text{major}}{(\mathrm y)}\ \xrightarrow[(2)\,\Delta]{(1)\,\mathrm{HNO_2+HBF_4}}\ \underset{\text{major}}{(\mathrm z)}$
\par Select the correct statement(s)
\begin{qoptions}
\opt{A}{(x) is \qfigi{p2_chem_s1_q03_2}}
\opt{B}{The product (y) is gives carbayl amine test}
\opt{C}{(z) is \qfigi{p2_chem_s1_q03_3}}
\opt{D}{conversion of (x) into (y) involves nucleophilic substitution}
\end{qoptions}
\qn{4} \qfigi{p2_chem_s1_q03_4}
\par Select the correct statement(s)
\begin{qoptions}
\opt{A}{The major product formed is an ester}
\opt{B}{The given reaction involves favorskii rearrangement}
\opt{C}{Carbocatic intermediate is involved}
\opt{D}{The major product is an ether}
\end{qoptions}
\qn{5} An aqueous solution of $6.3$ gm of a hydrated oxalic acid ($\mathrm{H_2C_2O_4\cdot xH_2O}$) is made up to $250$ ml. The $40$ ml of $0.1$ N NaOH was required to completely neutralize $10$ ml of the above prepared stock solution. Which of the following statement(s) are correct ?
\begin{qoptions}
\opt{A}{The acid is dihydrated}
\opt{B}{Equivalent weight of the hydrated salt is $45$}
\opt{C}{Equivalent weight of anhydrous acid is $45$}
\opt{D}{$20$ ml of the same stock solution would require $40$ ml $0.1$ M $\mathrm{Ca(OH)_2}$ for complete neutralisation}
\end{qoptions}
\qn{6} Which of the following statement is correct ?
\begin{qoptions}
\opt{A}{For any cyclic process, $\Delta U=0$}
\opt{B}{$\left(\dfrac{\delta U}{\delta V}\right)_T=\dfrac a{V^2}$ for $1$ mole of a gas obeying van der Waal's equation}
\opt{C}{$\overline C_p=\infty$ for water when it is in equilibrium with vapour at $1$ atm pressure and $373$ K.}
\opt{D}{For any real gas, the molar internal energy can be given as $\dfrac32R$}
\end{qoptions}
\qn{7} Which of the following reagent(s) would be most suitable to convert $\mathrm{(CH_3)_3N}$ into $\mathrm{(CH_3)_2NH}$
\begin{qoptions}
\opt{A}{HCl only}
\opt{B}{NCBr}
\opt{C}{$\mathrm{H_3O^+;\Delta}$}
\opt{D}{$\mathrm{(PPh_3+H_2O)}$ only}
\end{qoptions}
\qn{8} Consider the reaction given below :-
\par\noindent\qfig{p2_chem_s1_q08_1}
Select the correct statement(s)
\begin{qoptions}
\opt{A}{(X) can be \qfigi{p2_chem_s1_q08_2}}
\opt{B}{(X) can be \qfigi{p2_chem_s1_q08_3}}
\opt{C}{The reaction involves $\mathrm{E1_{CB}}$}
\opt{D}{The reaction involves E2 mechanism}
\end{qoptions}
\sect{Section-I(ii)}
\qn{9} \textbf{List -- I} (Metals) and \textbf{List -- II} (Method of extration)
\begin{qoptions}
\opt{P}{Fe}
\opt{Q}{Cu}
\opt{R}{Mg}
\opt{S}{Au}
\end{qoptions}
\begin{qoptions}
\opt{1}{Self reduction}
\opt{2}{Carbon reduction}
\opt{3}{Cyanide process}
\opt{4}{Electrolytic reduction}
\end{qoptions}
The correct option is -
\begin{qoptions}
\opt{A}{P$\to$2; Q$\to$1; R$\to$3; S$\to$4}
\opt{B}{P$\to$2; Q$\to$1; R$\to$4; S$\to$3}
\opt{C}{P$\to$1; Q$\to$3; R$\to$2; S$\to$4}
\opt{D}{P$\to$3; Q$\to$2; R$\to$4; S$\to$1}
\end{qoptions}
\qn{10} \textbf{List -- I}
\begin{qoptions}
\opt{P}{$\mathrm{PF_5}$}
\opt{Q}{$\mathrm{SF_6}$}
\opt{R}{$\mathrm{IF_7}$}
\opt{S}{$\mathrm{XeF_8^{2-}}$}
\end{qoptions}
\par\noindent\textbf{List -- II}
\begin{qoptions}
\opt{1}{Equatiorial bond longer than axial bond}
\opt{2}{Axial bond longer than equatorial bond}
\opt{3}{All bonds are of equal length}
\opt{4}{s-orbital is stereochemically inactive}
\end{qoptions}
The correct option is -
\begin{qoptions}
\opt{A}{P$\to$2; Q$\to$3; R$\to$1; S$\to$4}
\opt{B}{P$\to$2; Q$\to$3; R$\to$4; S$\to$1}
\opt{C}{P$\to$3; Q$\to$2; R$\to$1; S$\to$2}
\opt{D}{P$\to$3; Q$\to$1; R$\to$2; S$\to$4}
\end{qoptions}
\qn{11} Two mole of an ideal gas undergoes a cyclic $ABCA$ as shown in the diagram below. However, it is known that during the process $AB$, the gas uses only $25\%$ of energy absorbed as heat, in doing work.
\qfig{p2_chem_s2_q11_1}
\par\noindent\textbf{List -- I}
\begin{qoptions}
\opt{P}{Magnitude of $\Delta U$ for process $AB$}
\opt{Q}{Magnitude of $q$ for process $BC$}
\opt{R}{Magnitude of $\Delta H$ for process $CA$}
\opt{S}{Magnitude of total work done during cycle}
\end{qoptions}
\par\noindent\textbf{List -- II}
\begin{qoptions}
\opt{1}{$9.6$ KCal}
\opt{2}{$4.43$ KCal}
\opt{3}{$2.03$ KCal}
\opt{4}{$7.2$ Kcal}
\end{qoptions}
The correct option is -
\begin{qoptions}
\opt{A}{P$\to$2; Q$\to$1; R$\to$3; S$\to$4}
\opt{B}{P$\to$2; Q$\to$4; R$\to$1; S$\to$3}
\opt{C}{P$\to$4; Q$\to$2; R$\to$1; S$\to$3}
\opt{D}{P$\to$4; Q$\to$1; R$\to$2; S$\to$3}
\end{qoptions}
\qn{12} Given the following data\\
$\Delta_fH^\circ(\mathrm{CH_4,g})=-80\ \mathrm{kJ\,mol^{-1}}$\\
$\Delta_fH^\circ(\mathrm{C_2H_6,g})=-85\ \mathrm{kJ\,mol^{-1}}$\\
$\Delta_fH^\circ(\mathrm{C_2H_4,g})=50\ \mathrm{kJ\,mol^{-1}}$\\
Enthalpy of sublimation of Carbon (graphite) $=720\ \mathrm{kJ\,mol^{-1}}$\\
Enthalpy of atomisation of H $=215\ \mathrm{kJ\,mol^{-1}}$
\par\noindent\textbf{List -- I} (Thermodynamic data)
\begin{qoptions}
\opt{P}{Bond enthalpy of C--H bond (per mole)}
\opt{Q}{Bond enthalpy of C--C bond (per mole)}
\opt{R}{Bond enthalpy of C=C bond (per mole)}
\opt{S}{$\Delta H^\circ$, for the reaction : $n\mathrm{C_2H_4(g)}\longrightarrow\left(\mathrm{-CH-CH_2-}\right)_n\mathrm{(g)}$ (for $10$ moles of ethene)}
\end{qoptions}
\par\noindent\textbf{List -- II} (Values with respect to magnitude)
\begin{qoptions}
\opt{1}{$325$ kJ}
\opt{2}{$600$ kJ}
\opt{3}{$415$ kJ}
\opt{4}{$590$ kJ}
\end{qoptions}
The correct option is -
\begin{qoptions}
\opt{A}{P$\to$3; Q$\to$4; R$\to$1; S$\to$2}
\opt{B}{P$\to$3; Q$\to$4; R$\to$2; S$\to$1}
\opt{C}{P$\to$1; Q$\to$3; R$\to$2; S$\to$4}
\opt{D}{P$\to$3; Q$\to$1; R$\to$4; S$\to$2}
\end{qoptions}
\sect{Section-II}
\qn{1} Equivalent conductance of $0.2$M aqueous solution of a weak monobasic acid $HA$ is $10\ \mathrm{cm^2eq^{-1}}$ and that at infinite silution is $200\ \mathrm{cm^2\,eq^{-1}}$. What is the concentration of $\mathrm{H^+}$ ions.
\qn{2} A definite amount of solid $\mathrm{NH_4HS}$ is placed in a flask already containing ammonia gas at a fixed temperature and $0.1$ atm pressure. After the equilibrium is established, total pressure inside the flask is $1.1$ atm. Value of $K_p$ in $(\mathrm{atm})^2$ is :
\qn{3} Density of $\mathrm{X_{(g)}}$ at $1$ atm pressure and $500$ K is $0.8\ \mathrm{Kg/m^3}$. It effuses through a small hole at a rate of $4/5$ times that of $\mathrm O_2$ and a same condition find the compressibility factor of vapour.
\qn{4} Consider the given complex ion $\mathrm{\left[Pt(bn)_2\right]^{2+}}$, where $\mathrm{bn}$ stands for butylene diamine. Find the value of $\dfrac n2$ where $n$ is the total number of stereoisomers.
\qn{5} Find the value of $\dfrac n4$, where $n$ is the number of compounds will produce $\mathrm{POCl_3}$ when treated with $\mathrm{PCl_5}$ as intermediate product or final product.
\par\noindent $\mathrm{CH_3COOH}$, $\mathrm{P_4O_{10}}$, $\mathrm{H_2O}$, $\mathrm{H_2SO_4}$, $\mathrm{C_2H_5OH}$
\qn{6} Calculate total number of optically active stereoisomers for following compound?
\qfig{p2_chem_s3_q06_1}
\qn{7} Calculate the specific rotation of $(+)$ -- glyceraldehyde, if a solution containing $2.0$ gm/ $10$ml is placed in a $1$ gm polarimeter tube and its observed rotation at $25^\circ$C using sodium D line is $+1.74^\circ$.
\qn{8} Calculate the isoelectric point ($p$I) of lysine. Give the $\mathrm{pK_a}$ of $\alpha$-$\mathrm{NH_3}$ is $8.95$, $\mathrm{pK_a}$ side chain $\mathrm{NH_3}$ is $10.53$ and $\mathrm{pK_a}$ of $\alpha$-COOH is $2.18$.
\partsec{3}{MATHEMATICS}
\sect{Section-I(i)}
\qn{1} Let $O$ be the circumcentre of acute angled $\triangle ABC$ and let $r_1,r_2$ and $r_3$ be the radii of the circles drawn on the altitudes $OD$, $OE$ and $OF$ of triangle $OBC$, $OCA$ and $OAB$ as diameter, then minimum value of $\dfrac{a^2}{r_1^2}+\dfrac{b^2}{r_2^2}+\dfrac{c^2}{r_3^2}$ is divisible by :
\begin{qoptions}
\opt{A}{$36$}
\opt{B}{$72$}
\opt{C}{$144$}
\opt{D}{$216$}
\end{qoptions}
\qn{2} $f:[0,1]\to R$ is a function such that $f'(x)$ is continuous and $\displaystyle\int_0^1f(x)\,dx=0$, if maximum value of $|f'(x)|$ when $x\in[0,1]$ is $24$, the possible value of $\left|\displaystyle\int_0^xf(t)\,dt\right|\ x\in[0,1]$ is
\begin{qoptions}
\opt{A}{$1$}
\opt{B}{$2$}
\opt{C}{$3$}
\opt{D}{$4$}
\end{qoptions}
\qn{3} Let $A$ and $B$ are non singular matrices of order three, each of them commute with matrix $X$, where $AX=B$ and $XB=A$. If $X$ is not an identity matrix, then $\det(X^4+X^2+I)$ is divisible by (where $I$ is identity matrix of order three)
\begin{qoptions}
\opt{A}{$36$}
\opt{B}{$27$}
\opt{C}{$9$}
\opt{D}{$3$}
\end{qoptions}
\qn{4} Suppose that two complex number $z=a+ib$; $w=c+id$ satisfy the equation $\dfrac{z+w}{z}=\dfrac{w}{z+w}$, then (where $a,b,c,d\in R$ and $\sqrt{-1}=\mathrm i$)
\begin{qoptions}
\opt{A}{$\arg z-\arg w=\dfrac{2\pi}3$}
\opt{B}{both $b$ and $d$ are zero}
\opt{C}{$\arg z-\arg w=-\dfrac{2\pi}3$}
\opt{D}{at least one of $b$ and $d$ is non-zero}
\end{qoptions}
\qn{5} Let $a_n=\dfrac12+\dfrac13+\cdots+\dfrac1n$. Then
\begin{qoptions}
\opt{A}{$a_n$ is less than $\displaystyle\int_1^n\dfrac{dx}x$}
\opt{B}{$a_n$ is greater than $\displaystyle\int_2^n\dfrac{dx}x$}
\opt{C}{$\displaystyle\lim_{n\to\infty}\dfrac{a_n}{\ln n}=1$}
\opt{D}{$\displaystyle\lim_{n\to\infty}a_n$ is finite}
\end{qoptions}
\qn{6} Let $\alpha(x)$ and $\beta(x)$ be differentiable non linear functions on $[0,\infty)$ such that $\alpha(0)=1$, $\beta(0)=3$, $\beta(x)>2\alpha(x)$, $2\alpha(x)+\beta'(x)=5$ and $2\alpha'(x)+\beta(x)=5$ for all $x\in[0,\infty)$, then
\begin{qoptions}
\opt{A}{$\alpha(x)$ has one maxima in $[0,\infty)$}
\opt{B}{$\alpha(x)$ has one minima in $[1,\infty)$}
\opt{C}{$\beta(x)$ is increasing function in $[0,\infty)$}
\opt{D}{$\beta(x)$ is increasing function in $[1,\infty)$}
\end{qoptions}
\qn{7} Let $G$ be the centroid of the $\triangle ABC$ whose sides are of length $a$, $b$, $c$. If $P$ be a point in the plane of $\triangle ABC$ such that $PA=1$, $PB=3$, $PC=4$ and $PG=2$, then the value of $a^2+b^2+c^2$ is less than :
\begin{qoptions}
\opt{A}{$40$}
\opt{B}{$44$}
\opt{C}{$46$}
\opt{D}{$48$}
\end{qoptions}
\qn{8} If $y_n$ is $n$th derivative of $y$ w.r.t. $x$ then
\begin{qoptions}
\opt{A}{$y=(ax+b)^m$ then $y_n={}^mP_n\,a^n(ax+b)^{m-n}$}
\opt{B}{$y=\log(ax+b)$ then $y_n=(-1)^{n-1}\dfrac{(n-1)!\,a^n}{(ax+b)^n}$}
\opt{C}{$y=\sin(ax+b)$ then $y_n=a^n\sin\left(ax+b+n\dfrac\pi2\right)$}
\opt{D}{$y=\cos(ax+b)$ then $y_n=a^n\cos\left(ax+bx-n\dfrac\pi2\right)$}
\end{qoptions}
\sect{Section-I(ii)}
\qn{9} Match List -- I with List -- II and select the correct answer using the code given below the list.
\par\noindent\textbf{List -- I}
\begin{qoptions}
\opt{P}{$\displaystyle\sum_{r=1}^n\dfrac{2r+1}{\left(r^2+r\right)^2}$}
\opt{Q}{$\displaystyle\lim_{n\to\infty}\sum_{r=1}^n\dfrac{(2r+1)}{\left(r^2+r\right)^2}$}
\opt{R}{$\displaystyle\prod_{k=1}^n\left(\sum_{r=1}^k\dfrac{2r+1}{\left(r^2+r\right)^2}\right)$}
\opt{S}{$\displaystyle\lim_{n\to\infty}\prod_{k=1}^n\left(\sum_{r=1}^k\dfrac{2r+1}{\left(r^2+r\right)^2}\right)$}
\end{qoptions}
\par\noindent\textbf{List -- II}
\begin{qoptions}
\opt{1}{$\dfrac{n+2}{2(n+1)}$}
\opt{2}{$\dfrac{n(n+2)}{(n+1)^2}$}
\opt{3}{$1$}
\opt{4}{$\dfrac12$}
\end{qoptions}
\begin{qoptions}
\opt{A}{P$\to$2; Q$\to$3; R$\to$1; S$\to$4}
\opt{B}{P$\to$3; Q$\to$1; R$\to$2; S$\to$4}
\opt{C}{P$\to$1; Q$\to$3; R$\to$2; S$\to$4}
\opt{D}{P$\to$4; Q$\to$1; R$\to$2; S$\to$3}
\end{qoptions}
\qn{10} Match List-I with List-II and select the correct answer using the code given below the list. If $y+x=4$ is a tangent to the hyperbola at the point $(2,2)$ intersects asymptotes of hyperbola at $A$ and $B$. If length of $AB$ is $4\sqrt2$ and centre of hyperbola is at origin, then match the following List--I with List--II
\par\noindent\textbf{List--I}
\begin{qoptions}
\opt{P}{The equation of conjugate axis of hyperbola is}
\opt{Q}{The equation of hyperbola is}
\opt{R}{The equation of directrix can be}
\opt{S}{The equation of transverse axis is}
\end{qoptions}
\par\noindent\textbf{List--II}
\begin{qoptions}
\opt{1}{$x+y=2\sqrt2$}
\opt{2}{$y=x$}
\opt{3}{$y=-x$}
\opt{4}{$xy=4$}
\end{qoptions}
\begin{qoptions}
\opt{A}{P$\to$4; Q$\to$3; R$\to$1; S$\to$2}
\opt{B}{P$\to$3; Q$\to$4; R$\to$1; S$\to$2}
\opt{C}{P$\to$3; Q$\to$1; R$\to$4; S$\to$2}
\opt{D}{P$\to$3; Q$\to$4; R$\to$2; S$\to$1}
\end{qoptions}
\qn{11} Consider four distinct circles with radii $3$, $r_1$, $r_2$, $r_3$ ($r_1,r_2,r_3>3$) touches all three fixed distinct lines in a plane, where $r_1,r_2,r_3$ are roots of equation $x^3-mx^2+243x-n=0$, $m,n\in R$
\par\noindent\textbf{List-I}
\begin{qoptions}
\opt{P}{The value of $m$ is}
\opt{Q}{The value of $r_1$ is}
\opt{R}{The value of $r_1+r_2$ is}
\opt{S}{The value of $n$ is}
\end{qoptions}
\par\noindent\textbf{List-II}
\begin{qoptions}
\opt{1}{$729$}
\opt{2}{$18$}
\opt{3}{$9$}
\opt{4}{$27$}
\end{qoptions}
\begin{qoptions}
\opt{A}{P$\to$3; Q$\to$4; R$\to$2; S$\to$1}
\opt{B}{P$\to$4; Q$\to$3; R$\to$1; S$\to$2}
\opt{C}{P$\to$4; Q$\to$3; R$\to$2; S$\to$1}
\opt{D}{P$\to$1; Q$\to$4; R$\to$2; S$\to$1}
\end{qoptions}
\qn{12} Given $f:A\to B$.
\qfig{p2_math_s2_q12_1}
In the List-I there are some conditions and in the List-II there are number of functions satisfying the conditions in the List-II.\\
Match List-I with List-II and select the correct answer using the code given below the list.
\par\noindent\textbf{List-I}
\begin{qoptions}
\opt{P}{$f(1)>f(2)>f(3)>f(4)$}
\opt{Q}{$f(i)\neq f(j)$ for $i\neq j$ (where $i,j\in A$)}
\opt{R}{$f(1)\ge f(2)\ge f(3)\ge f(4)$}
\opt{S}{$f(1)<f(2)<f(3)>f(4)$ and $f(i)\neq f(j)$ for $i\neq j$ (where $i,j\in A$)}
\end{qoptions}
\par\noindent\textbf{List-II}
\begin{qoptions}
\opt{1}{${}^{10}C_4\cdot{}^3C_1$}
\opt{2}{${}^{10}C_4\cdot24$}
\opt{3}{${}^{10}C_4$}
\opt{4}{${}^{13}C_4$}
\end{qoptions}
\begin{qoptions}
\opt{A}{P$\to$2; Q$\to$3; R$\to$4; S$\to$1}
\opt{B}{P$\to$3; Q$\to$2; R$\to$4; S$\to$1}
\opt{C}{P$\to$3; Q$\to$1; R$\to$4; S$\to$2}
\opt{D}{P$\to$1; Q$\to$4; R$\to$3; S$\to$2}
\end{qoptions}
\sect{Section-II}
\qn{1} Two lines are drawn at right angle one being a tangent to $y^2=12x$ and the other to $x^2=16y$. If the locus of their point of intersection is $(x^2+y^2)(\ell x+my)+(nx-3y)^2=0$ then $\ell+n-m$ is equal to
\qn{2} Let $AP(a,d)$ denote the set of all the terms of an arithmetic progression containing $2020$ terms with first term $a$ and common difference $d$. If $AP(1,3)$ and $AP(9,7)$ have $\lambda$ number of uncommon terms, then the value of $\dfrac\lambda{100}$ is
\qn{3} If $P$ lies in the first quadrant on the ellipse $\dfrac{x^2}{a^2}+\dfrac{y^2}{b^2}=1$ (where $a>b$), and tangent \& normal drawn at $P$ meets major axis at the points $T$ \& $N$ respectively, then the value of $\dfrac{4\left(|F_2N|+|F_1N|\right)\left(|F_2T|-|F_1T|\right)}{5\left(|F_2N|-|F_1N|\right)\left(|F_2T|+|F_1T|\right)}$ is equal to\\
(where $F_1$ \& $F_2$ are the foci $(ae,0)$ \& $(-ae,0)$ respectively)
\qn{4} Let $f:\mathbb R\to\mathbb R$ be a polynomial function with real coefficients such that $(x^2+x+1)f(x-1)=(x^2-x+1)f(x)\ \forall\ x\in\mathbb R$ and $f(0)=1$, then the value of $\dfrac{f(1)}{10}$ is
\qn{5} Let $\operatorname{diag}.(a,b,c)$ represents a $3\times3$ diagonal matrix with diagonal elements $a,b,c$ (in order). Now, consider matrices $A=\operatorname{diag}.(3,2,5)$, $B=\operatorname{diag}.(11,19,13)$ and $C=\operatorname{diag}.(61,53,31)$. The number of factors of trace of the matrix `$ABC$' is\\
(trace $=$ sum of diagonal elements)
\qn{6} If $a,b,c$ are positive real numbers and the minimum value of $\dfrac{a+3c}{a+c+2b}+\dfrac{4b}{a+b+2c}-\dfrac{8c}{a+b+3c}$ is $m$, then the value of $|m|$ is $(\text{use }\sqrt2=1.4)$
\qn{7} In a tetrahedron $OABC$, the edges are of lengths, $|OA|=|BC|=a$, $|OB|=|AC|=b$, $|OC|=|AB|=c$. Let $G_1$ and $G_2$ be the centroids of the triangle $ABC$ and $AOC$ such that $OG_1\perp BG_2$, then the value of $\dfrac{a^2+c^2}{b^2}$ is
\qn{8} A game of luck is designed in which you have to shoot a target. The organizer puts two bullet in an empty six round revolver and spins. You point it towards balloon and shoot. It was an empty shot (i.e. bullet was not fired). If the bullets are placed randomly then the probability that you will shoot the balloon in next shot is $\dfrac PQ$ (Where $P$ and $Q$ are relatively coprimes) \{It is assumed that fired bullet always hits the target\}, then the value of $(P^2+Q^2)$ is
\end{multicols}
\end{document}
